% !TeX program = pdflatex
% Companion (applied) paper to Part VI of the godsil-gutman-lean series.
% The Cycles in the Chip — a machine-checked short-cycle census of the IEEE 802.11n
% LDPC codes via the trace-formula gap law.
% Build: pdflatex ldpc-census.tex  (twice, for refs)
% Data: research/ldpc-gaplaw/ (CENSUS.md, census_out.txt, census_pilot2.py).
% Theory it leans on: Part VI (gap-window-lean), gap law sorry-free on 1<=k<=g+1.

\documentclass[11pt]{article}
\usepackage[T1]{fontenc}
\usepackage{lmodern}
\usepackage{amsmath,amssymb,amsthm,booktabs,graphicx,hyperref}
\usepackage{xurl}
\usepackage[table]{xcolor}
\definecolor{ggteal}{HTML}{1B6F8C}
\definecolor{ggband}{HTML}{DCE7EB}
\definecolor{ggamber}{HTML}{E08A1E}
\usepackage[margin=1in]{geometry}
\usepackage{microtype}
\usepackage[most]{tcolorbox}
\usepackage{tikz}
\usetikzlibrary{arrows.meta,positioning}
\setlength{\emergencystretch}{2em}
\newtheorem{theorem}{Theorem}
\newcommand{\lean}[1]{\texttt{#1}}
\newcommand{\Z}{\mathbb{Z}}
\newcommand{\C}{\mathbb{C}}
\DeclareMathOperator{\tr}{tr}
\hypersetup{colorlinks=true,linkcolor=ggteal,citecolor=ggteal,urlcolor=ggteal}
\newtcbox{\keybox}{on line,boxrule=0.9pt,colframe=ggteal,colback=ggband!40,
  arc=3pt,left=8pt,right=8pt,top=5pt,bottom=5pt}
\newcommand{\keyeq}[1]{\begin{center}\keybox{$\displaystyle #1$}\end{center}}

\title{\bf Certified Short-Cycle Counts for the IEEE 802.11n (WiFi) LDPC Codes\\[4pt]
  \large\normalfont A machine-checked census of the deployed parity-check graphs\\
  via the trace-formula gap law}
\author{Carles Mar\'in\thanks{Census computed with AI assistance (Claude, Anthropic); the
mathematics, the theorem it rests on, and all claims are the author's responsibility. The
certified identity is proved in the companion paper~\cite{MarinGapWindow2026}.}\\[3pt]
  \normalsize Independent researcher\\
  \normalsize \texttt{karlesmarin@gmail.com}}
\date{June 2026}

\begin{document}
\maketitle

\begin{abstract}
Short cycles in the Tanner graph of a low-density parity-check (LDPC) code degrade
belief-propagation decoding, and counting them is a standard task in code design. This paper
does not count them faster. It counts them \emph{with a certificate}: for each of the four
LDPC codes of the IEEE 802.11n (WiFi) standard at block length $n = 648$, the number of
shortest cycles is obtained from a machine-checked identity---the trace-formula gap law,
proved \lean{sorry}-free in Lean~4 in the companion paper~\cite{MarinGapWindow2026}---and
cross-checked by three mutually independent computations that agree exactly: the
non-backtracking trace $\tr B^g/2g$, the gap law $(\tr A^g - p_g)/2g$, and direct
enumeration. The census reads $c_6 = 3942$ (rate $1/2$), $c_6 = 8046$ (rate $2/3$),
$c_4 = 54$ (rate $3/4$, where the girth collapses to $4$) and $c_6 = 32\,346$ (rate $5/6$),
and it tells a design story: the rate-$3/4$ code pays for its density with four-cycles, the
most damaging kind, while the rate-$5/6$ code buys its girth back at the price of thirty-two
thousand hexagons. The counting algorithms of the coding-theory literature are faster and
reach further; the contribution here is the certification layer---an exact identity backed
by a theorem with a checkable proof---which, to the best of my knowledge, no prior
cycle-counting tool carries.
\end{abstract}

\noindent\emph{Reader's map.} The theorem is borrowed (Section~\ref{sec:identity}; proved in
the companion). The new thing is the audited census: three independent routes
(Section~\ref{sec:routes}), the four deployed codes (Section~\ref{sec:census},
Table~\ref{tab:census}), and the honest accounting of what the certificate does and does not
buy (Sections~\ref{sec:buys}--\ref{sec:boundary}).

% ---------------------------------------------------------------
\section{Introduction: a number a decoder lives by}
\label{sec:intro}

What does it mean to \emph{trust} a cycle count? When Gallager invented LDPC codes in 1962
the world had no use for them; thirty years later they returned---to WiFi, to 5G, to deep
space---decoded by belief propagation on a bipartite \emph{Tanner graph}~\cite{Tanner1981}
whose two sides are the code's bits and its parity checks~\cite{Gallager1962,MacKay1999}.
Belief propagation is exact on a tree and only approximate on a graph with cycles: a short
cycle lets a message circle back to its origin and reinforce itself, a rumor mistaken for
independent confirmation. The shorter the cycle, the tighter the loop, so the \emph{girth}
of the Tanner graph---the length of its shortest cycle---and the \emph{number} of cycles at
and just above the girth became quantities a code designer measures~\cite{Tanner1981,
KarimiBanihashemi2013}.

Measuring them is a solved problem. Spectral and message-passing algorithms count short
cycles efficiently, and on the quasi-cyclic graphs of real standards they count well past the
girth~\cite{KarimiBanihashemi2013,BlakeLin2018,Dehghan2018}. This paper adds nothing to their
speed or reach. It adds one thing they lack: a count whose defining identity is a
\emph{theorem with a machine-checked proof}. The companion paper~\cite{MarinGapWindow2026}
proves, \lean{sorry}-free in Lean~4 over Mathlib, that for a finite graph of girth $g$ the
number of shortest cycles is
\[
c_g \;=\; \frac{\tr A^g - p_g}{2g} \;=\; \frac{\tr B^g}{2g},
\]
with $A$ the adjacency matrix, $B$ the Hashimoto non-backtracking operator, and $p_g$ the
$g$-th power sum of the roots of the matching polynomial. Here we apply that identity, as a
\emph{certified census}, to the LDPC codes a WiFi chip actually runs.

% ---------------------------------------------------------------
\section{The identity we apply, and how far it is certified}
\label{sec:identity}

Which theorem are we leaning on, and where does its guarantee stop? The companion
paper~\cite{MarinGapWindow2026} proves the \emph{trace-formula gap law}: for every $k$ on the
sharp window $1 \le k \le g+1$,
\keyeq{\tr A^k \;-\; p_k \;=\; \tr B^k .}
Below the girth ($k < g$) both sides vanish; at $k \in \{g, g+1\}$ both equal $2k\,c_k$,
where $c_k$ is the number of $k$-cycles. Two consequences are all we need. First, the
below-girth vanishing $\tr B^k = 0$ for $k < g$ is a sanity certificate: a Tanner graph of
girth $g$ \emph{must} report zero closed non-backtracking walks shorter than $g$, and any code
graph failing it is mis-parsed. Second, at $k = g$ the law is the shortest-cycle count above.

\emph{How far is it certified?} Both facts are theorems in Lean, depending only on the three
standard axioms (\lean{propext}, \lean{Classical.choice}, \lean{Quot.sound}); the
below-girth vanishing is ``Stone~A'' and the window equality is the capstone ``Stone~B'' of
the companion. So the \emph{identity} that turns spectra and matching counts into a cycle
count is machine-checked end to end. What is \emph{not} part of the theorem is the arithmetic
on any particular matrix: evaluating $\tr A^g$, $p_g$ or $\tr B^g$ on a concrete $2376$-edge
Tanner graph is ordinary exact computation, not a Lean proof. That gap---between a certified
formula and its evaluation on a specific input---is exactly what the three-route cross-check
of the next section is built to close.

% ---------------------------------------------------------------
\section{Three routes to one number}
\label{sec:routes}

Why compute the same integer three ways? Because a certified \emph{formula} evaluated by
fallible code is only as trustworthy as the evaluation, and the cheapest guard against an
implementation slip is redundancy by independent method. The census computes each $c_g$ along
three routes that share no code path and meet only at the answer.

\begin{tcolorbox}[colframe=ggteal!70,colback=ggband!20,arc=3pt,boxrule=0.8pt,
  title=\textbf{The three routes (mutually independent)}, coltitle=white,colbacktitle=ggteal]
\textbf{R1 --- non-backtracking trace.} Build the Hashimoto operator $B$ on the $2|E|$ darts
as a sparse $0/1$ matrix and form $\tr B^g$ without ever densifying (powers by sparse
products, traces by the Hadamard identity $\tr B^k = \sum (B^a \circ (B^b)^{\!\top})$,
$a+b=k$); then $c_g = \tr B^g / 2g$.\\[2pt]
\textbf{R2 --- the gap law.} Compute $\tr A^g$ from the adjacency matrix and $p_g$ by Newton
from matching counts---$p_4 = 2m_1^2 - 4m_2$, $p_6 = 2m_1^3 - 6m_1m_2 + 6m_3$, with $m_1 =
|E|$, $m_2 = \binom{m_1}{2} - \sum_v \binom{d_v}{2}$, and $m_3$ by an exact bipartite
edge-deletion recursion; then $c_g = (\tr A^g - p_g)/2g$.\\[2pt]
\textbf{R3 --- direct enumeration.} Count cycles combinatorially: for $g = 4$, sum
$\binom{\text{codeg}}{2}$ over same-side vertex pairs (bipartite four-cycles); for $g = 6$, a
VF2 subgraph search for $C_6$, divided by its $12$ automorphic rootings.
\end{tcolorbox}

R1 lives in the geodesic world of the non-backtracking operator; R2 in the spectral-plus-
matching world of the gap law; R3 in plain enumeration. They invoke disjoint libraries and
disjoint arithmetic. When all three return the same integer, an error would have to be a
\emph{coincidence across three methods}---and the gap law is exactly the theorem that says R1
and R2 must agree, so a three-way match also re-confirms the identity on each real graph.

The matching-count route deserves one remark. The recursion $m_3 = \tfrac13 \sum_{(u,v)\in E}
m_2(G - \{u,v\})$ is evaluated in closed form on the bipartite graph---$N(u)$ and $N(v)$ are
disjoint, so the deletion's matching numbers reduce to degree sums---never copying the graph;
this is what keeps R2 a few seconds rather than a few hours on a $2376$-edge instance.

% ---------------------------------------------------------------
\section{The census: four codes a WiFi chip runs}
\label{sec:census}

What do the deployed codes actually look like? IEEE 802.11n specifies LDPC codes at three
block lengths; we take the longest, $n = 648$, with quasi-cyclic lifting size $z = 27$, in
all four code rates. Each parity-check matrix yields a Tanner graph on $|V| = n + m$ vertices
($m = n(1-\text{rate})$ checks) with a constant $|E| = 2376$ edges. The census, with the
below-girth zeros confirmed on every code, is Table~\ref{tab:census}.

\begin{table}[h]
\centering\small
\rowcolors{2}{ggband!35}{white}
\begin{tabular}{@{}lrrrrl@{}}
\toprule
\rowcolor{ggteal!18} rate & $|V|$ & $|E|$ & girth $g$ & $c_g$ (R1\,=\,R2\,=\,R3) & shortest cycle \\
\midrule
$1/2$ & 972 & 2376 & 6 & $3\,942$  & hexagon \\
$2/3$ & 864 & 2376 & 6 & $8\,046$  & hexagon \\
$3/4$ & 810 & 2376 & 4 & $54$      & \textbf{square} \\
$5/6$ & 756 & 2376 & 6 & $32\,346$ & hexagon \\
\bottomrule
\end{tabular}
\caption{The certified short-cycle census of the IEEE 802.11n LDPC codes, $n = 648$,
$z = 27$. For each code the three independent routes of Section~\ref{sec:routes} return the
same integer, and the Lean-certified below-girth vanishing $\tr B^k = 0$ ($k < g$) holds on
all four Tanner graphs. Computed in exact arithmetic; the pipeline is
\lean{research/ldpc-gaplaw/census\_pilot2.py}.}
\label{tab:census}
\end{table}

\emph{The story the four numbers tell.} Read down the table and a design trade-off appears
(Figure~\ref{fig:tradeoff}). At rate $1/2$ the graph is sparse enough to hold girth $6$ with
the fewest hexagons, $3942$---the cleanest of the four. At rate $2/3$ the girth survives but
the hexagon population doubles. At rate $3/4$ the girth \emph{collapses} to $4$: the density
forces $54$ squares, the most damaging cycles a Tanner graph can carry. At rate $5/6$, denser
still, the designers recover girth $6$---trading the squares away---but pay for it with
$32\,346$ hexagons, an eight-fold jump over rate $1/2$. The certificate does not judge the
trade; it makes the price exact.

\begin{figure}[t]
\centering
\begin{tikzpicture}[xscale=2.4,yscale=1.05]
  % axes
  \draw[->,ggteal!70,semithick] (-0.2,0) -- (4.1,0) node[right,font=\footnotesize]{rate};
  \draw[->,ggteal!70,semithick] (-0.2,0) -- (-0.2,5.3)
     node[above,font=\footnotesize,align=center]{$\log_{10} c_g$};
  \foreach \y/\l in {1.73/54, 3.60/3942, 3.91/8046, 4.51/32346}
    \draw[ggband] (-0.25,\y) -- (4.05,\y);
  % data points: x = rate index, y = log10(c_g); girth color
  % rate 1/2 -> x=0.5  c=3942  g6 ; 2/3 -> x=1.5 c=8046 g6 ; 3/4 -> x=2.5 c=54 g4 ; 5/6 -> x=3.5 c=32346 g6
  \node[circle,fill=ggteal,inner sep=1.7pt,label={[font=\scriptsize]above:$3942$}]  at (0.5,3.60){};
  \node[circle,fill=ggteal,inner sep=1.7pt,label={[font=\scriptsize]above:$8046$}]  at (1.5,3.91){};
  \node[circle,fill=ggamber,inner sep=2.1pt,label={[font=\scriptsize]below:$54$}]    at (2.5,1.73){};
  \node[circle,fill=ggteal,inner sep=1.7pt,label={[font=\scriptsize]above:$32346$}] at (3.5,4.51){};
  % girth labels under axis
  \node[font=\scriptsize] at (0.5,-0.45){$1/2$};
  \node[font=\scriptsize] at (1.5,-0.45){$2/3$};
  \node[font=\scriptsize] at (2.5,-0.45){$3/4$};
  \node[font=\scriptsize] at (3.5,-0.45){$5/6$};
  \node[font=\scriptsize,ggteal]  at (0.5,-0.85){$g=6$};
  \node[font=\scriptsize,ggteal]  at (1.5,-0.85){$g=6$};
  \node[font=\scriptsize,ggamber] at (2.5,-0.85){$g=4$};
  \node[font=\scriptsize,ggteal]  at (3.5,-0.85){$g=6$};
  % connect the girth-6 trend
  \draw[ggteal!55,dashed] (0.5,3.60) -- (1.5,3.91) -- (3.5,4.51);
  \node[ggamber,font=\scriptsize,align=center] at (2.5,2.5){the girth-$4$\\collapse};
\end{tikzpicture}
\caption{The design landscape of the four IEEE 802.11n codes at $n = 648$. The three
girth-$6$ rates trace a rising hexagon count (dashed); rate $3/4$ drops out of that trend
because its girth collapses to $4$, where the count is small ($54$) but the cycles are the
short, decoder-damaging kind. ``Few cycles'' and ``long cycles'' are different goods, and the
census separates them.}
\label{fig:tradeoff}
\end{figure}

% ---------------------------------------------------------------
\section{What the certificate buys}
\label{sec:buys}

Counting short cycles is solved---so what is new here? Not the count, and not its cost. The
spectral method of Blake and Lin and the message-passing methods of Karimi and Banihashemi
count further (to $2g-2$ and beyond, with correction terms) and faster, and they exploit the
quasi-cyclic block structure that the present pipeline ignores~\cite{BlakeLin2018,
KarimiBanihashemi2013,Dehghan2018}. Against them this census is slower and shorter-reaching.

What it carries that they do not is a \emph{certificate}: the identity producing the number
is a theorem with a proof a machine has checked, not a derivation trusted to a paper. Three
things follow that matter in practice. The below-girth vanishing is a \emph{free integrity
check} on a parsed parity-check matrix---a single mis-read entry that creates a phantom short
cycle shows up as a nonzero $\tr B^k$ where the theorem forbids one. The agreement of the
non-backtracking and matching routes is not a lucky coincidence but a \emph{predicted}
equality, so a mismatch localizes a bug to one route. And the smallest moments of the
non-backtracking operator, $\tr B^g = 2g\,c_g$ and $\tr B^{g+1} = 2(g{+}1)c_{g+1}$, which
govern percolation and epidemic thresholds on a network~\cite{KarrerNewman2014} and the
detectability limit in community detection~\cite{Krzakala2013}, are now pinned exactly from
quantities computed by independent, well-tested code---an entry-level regression test for any
implementation of $B$.

% ---------------------------------------------------------------
\section{Honest boundaries}
\label{sec:boundary}

Where exactly does the guarantee stop? In three places, stated plainly.

\emph{The window is $k \le g+1$.} The certified identity counts cycles \emph{at} the girth
(and at $g+1$); the engineering quantities past it---multiplicities up to $2g-2$, the
correction terms of~\cite{BlakeLin2018,Dehghan2018}---are not covered by the theorem and are
not computed here.

\emph{The matrices are third-party.} The Tanner graphs are built from parity-check matrices
derived from the IEEE 802.11n standard via a public toolkit, not re-derived from the
standard's Annex~F. Their quasi-cyclic block structure ($z = 27$) and degree profiles were
cross-checked, and the three-route agreement is itself strong evidence the graphs are the
intended ones, but a standards-grade census would re-derive $H$ from the normative document.
The certificate covers the \emph{counting identity}, not the provenance of the input.

\emph{No algorithmic novelty.} The routes are textbook. The contribution is their
\emph{composition under a machine-checked identity}, which a literature search on 2026-06-11
found absent: no Lean, Coq, or other proof-assistant work on cycle counting for coding theory.
As always, such a claim is made only to the best of one's knowledge.

% ---------------------------------------------------------------
\section{Open questions}
\label{sec:questions}

\begin{enumerate}
\item \emph{The 5G headline.} The same pipeline applies to the 3GPP NR codes (TS 38.212,
base graphs BG1 and BG2 with their deployed lifting sizes). A certified census across the
5G rate-compatible family is the natural next target---larger graphs, the same theorem.
\item \emph{The corrected window in Lean.} Engineers count to $2g-2$ with closed-form
corrections; the companion paper leaves the first correction term ($k = g+2$, the lollipop
walks) as an open formalization. A certified count that reaches $2g-2$ would match the
literature's reach \emph{with} a proof.
\item \emph{Re-derivation from the standard.} Building $H$ directly from IEEE 802.11n
Annex~F and 3GPP TS 38.212, inside the certified pipeline, would close the provenance gap and
make the census standards-grade end to end.
\item \emph{The certificate as a design instrument.} The pair $(\tr B^g, \tr B^{g+1})$ is a
cheap certified invariant of a Tanner graph. How much of a code's short-cycle profile---and
of its decoding behavior---does this certified pair pin down, and where does it stop
distinguishing codes that decode differently?
\end{enumerate}

\section*{Data and code availability}
The parity-check alist files, the census pipeline
(\lean{research/ldpc-gaplaw/census\_pilot2.py}) and the full output are in the public
repository \url{https://github.com/karlesmarin/godsil-gutman-lean}, alongside the Lean
sources of the companion paper~\cite{MarinGapWindow2026} that prove the identity applied here.

\begin{thebibliography}{99}\small

\bibitem{Gallager1962} R.~G. Gallager,
\emph{Low-density parity-check codes}, IRE Trans. Inform. Theory \textbf{8} (1962) 21--28.

\bibitem{Tanner1981} R.~M. Tanner,
\emph{A recursive approach to low complexity codes}, IEEE Trans. Inform. Theory \textbf{27}
(1981) 533--547.

\bibitem{MacKay1999} D.~J.~C. MacKay,
\emph{Good error-correcting codes based on very sparse matrices}, IEEE Trans. Inform. Theory
\textbf{45} (1999) 399--431.

\bibitem{KarimiBanihashemi2013} M. Karimi, A.~H. Banihashemi,
\emph{Message-passing algorithms for counting short cycles in a graph},
IEEE Trans. Commun. \textbf{61} (2013) 485--495.

\bibitem{BlakeLin2018} I.~F. Blake, S. Lin,
\emph{On short cycle enumeration in biregular bipartite graphs},
IEEE Trans. Inform. Theory \textbf{64} (2018) 6526--6535.

\bibitem{Dehghan2018} A. Dehghan, A.~H. Banihashemi,
\emph{On computing the multiplicity of cycles in bipartite graphs using the degree
distribution and the spectrum of the graph}, IEEE Trans. Inform. Theory \textbf{65} (2019)
3778--3789.

\bibitem{KarrerNewman2014} B. Karrer, M.~E.~J. Newman, L. Zdeborov\'a,
\emph{Percolation on sparse networks}, Phys. Rev. Lett. \textbf{113} (2014) 208702.

\bibitem{Krzakala2013} F. Krzakala, C. Moore, E. Mossel, J. Neeman, A. Sly, L. Zdeborov\'a,
P. Zhang, \emph{Spectral redemption in clustering sparse networks},
Proc. Natl. Acad. Sci. \textbf{110} (2013) 20935--20940.

\bibitem{Hashimoto1989} K. Hashimoto,
\emph{Zeta functions of finite graphs and representations of $p$-adic groups},
Adv. Stud. Pure Math. \textbf{15} (1989) 211--280.

\bibitem{Mathlib2020} The mathlib community,
\emph{The Lean mathematical library}, in: CPP 2020, ACM, 2020, 367--381.

\bibitem{MarinGapWindow2026} C. Mar\'in,
\emph{The walks that remember the cycles: a machine-checked sharp gap law between the
matching polynomial and the non-backtracking spectrum in Lean~4} (Part VI of the
godsil--gutman--lean series), Zenodo, 2026.

\end{thebibliography}

\end{document}
